\subsection{Injective envelopes}

DEFINITION 4. --- Let M be an A-module. An injective envelope of M is a pair (I, i), where I is an injective A-module and i : M → I a homomorphism possessing the following property:

\begin{enumerate}
\def\labelenumi{(\Alph{enumi})}
\setcounter{enumi}{4}
\tightlist
\item
  for a submodule P of I to be zero, it is necessary and sufficient that \(i^{-1}(\mathbf{P})\) be zero. Note that (E) implies that \(i\) is injective. One often identifies M with the submodule \(i(\mathbf{M})\) of I, and one then says that I is an injective envelope of M.
\end{enumerate}

Example 1. Suppose A is an integral domain and M is torsion-free. Let K be the field of fractions of A and \(i: \mathbf{M} \to \mathbf{K} \otimes_{\mathbf{A}} \mathbf{M}\) the canonical homomorphism. Then \((\mathbf{K} \otimes_{\mathbf{A}} \mathbf{M}, i)\) is an injective envelope of M (II, p.~116, prop. 26 and X, p.~17, cor. 2).

THEOREM 2. --- Let M be an A-module.

\begin{enumerate}
\def\labelenumi{\alph{enumi})}
\item
  M has injective envelopes.
\item
  If (I, i) and (J, j) are two injective envelopes of M, there exists an isomorphism \(f: I \to J\) such that \(f \circ i = j\) .
\end{enumerate}

Note that the homomorphism \(f\) whose existence is asserted in \(b\) ) is not uniquely determined in general.

\begin{enumerate}
\def\labelenumi{\alph{enumi})}
\tightlist
\item
  One may assume that M is a submodule of an injective A-module E (X, p.~19, cor. 3). Consider the set ordered by inclusion \(\mathcal{F}\) of submodules I of E, containing M, and such that the canonical injection \(i: \mathrm{M} \to \mathrm{I}\) satisfies property (E). Since \(\mathcal{F}\) is inductive, it has a maximal element (E, III, p.~20); let I be a maximal element of \(\mathcal{F}\) . It suffices to prove that I is a direct summand submodule of E. Let N be a submodule of E such that \(\mathrm{N} \cap \mathrm{I} = 0\) , and maximal for this property (such an N exists by loc. cit.). The composed homomorphism
\end{enumerate}

\[
\mathrm{I} \xrightarrow {u} \mathrm{E} \xrightarrow {v} \mathrm{E} / \mathrm{N},
\]

where \(u\) and \(v\) are the canonical homomorphisms, is injective; since \(E\) is injective, there therefore exists a homomorphism \(w: E/N \to E\) such that \(w \circ v \circ u = u\) , that is, \(w \circ v(x) = x\) for \(x \in I\) . Set \(I' = \text{Im}(w) = \text{Im}(w \circ v)\) and let \(i': M \to I'\) be the canonical injection. Then \(I \subset I'\) ; to complete the proof, it suffices to show that \(i'\) satisfies condition (E): this implies that \(I = I'\) (maximal character of I) , hence that \(w \circ v\) is a projector of \(E\) over \(I\) .

Let P then be a submodule of I' such that \(P \cap M = 0\) . We have \(P = w \circ v(Q)\) , where Q is a submodule of E containing N; moreover

\[
\mathbf {Q} \cap \mathbf {M} = w \circ v (\mathbf {Q} \cap \mathbf {M}) \subset \mathbf {P} \cap \mathbf {M} = 0,
\]

Thus Q ∩ I = 0 since i : M → I has property (E). By the maximality of N, this implies Q = N, that is, v(Q) = 0, hence P = 0, as was to be shown.

\begin{enumerate}
\def\labelenumi{\alph{enumi})}
\setcounter{enumi}{1}
\tightlist
\item
  Let (I, i) and (J, j) be two injective envelopes of M. Since J is injective, there exists a homomorphism \(f: I \to J\) such that \(f \circ i = j\) ; we have
\end{enumerate}

\[
i ^ {- 1} (\operatorname{Ker} f) = \operatorname{Ker} j = 0,
\]

hence Ker f = 0 and f is injective. Then \(f(I)\) is an injective submodule of J, hence a direct summand; since j satisfies (E), this implies \(f(I) = J\) and f is bijective.

Remarks. --- 1) Let (I, i) be an injective envelope of M and let j : M → J be an injective homomorphism from M into an injective A-module J. By the proof above, there exists an injective homomorphism f : I → J such that f ∘ i = j.

\begin{enumerate}
\def\labelenumi{\arabic{enumi})}
\setcounter{enumi}{1}
\tightlist
\item
  Let (I, i) be an injective envelope of M. Identify M with the submodule \(i(\mathbf{M})\) of I. For every submodule N of M, there exists an injective submodule of I which is an injective envelope of N (apply Remark 1 to N). Conversely, every injective submodule J of I is an injective envelope of \(J \cap M\) .
\end{enumerate}

PROPOSITION 14. --- Let I be a nonzero injective A-module. The following conditions are equivalent:

\begin{enumerate}
\def\labelenumi{(\roman{enumi})}
\item
  I is indecomposable (VII, § 4, no. 8, def. 3);
\item
  0 is not the intersection of two nonzero submodules of I;
\item
  I is the injective envelope of all of its nonzero submodules;
\item
  the ring \(\operatorname{End}_{\mathbf{A}}\) (I) is local (VIII, § 1, no. 4, def. 4).
\item
  \(\Rightarrow\) (iii): let M be a nonzero submodule of I. By Remark 2 above, there exists a submodule I' of I which is an injective envelope of M. Since I' is nonzero and a direct factor of I, we have I = I' if I is indecomposable.
\item
  \(\Rightarrow\) (ii): let E and F be two submodules of I, such that \(\mathbf{E} \cap \mathbf{F} = 0\) . If \(\mathbf{E} \neq 0\) , then I is an injective envelope of E by (iii), so " \(\mathbf{E} \cap \mathbf{F} = 0\) " implies " \(\mathbf{F} = 0\) ".
\item
  \(\Rightarrow\) (i): this is trivial.
\item
  \(\Rightarrow\) (i): this follows from VIII, § 1, no. 6, prop. 13.
\item
  \(\Rightarrow\) (iv): suppose I is indecomposable. Note first that every injective endomorphism \(f\) of I is bijective (since \(f(\mathbf{I})\) is then a nonzero direct summand of I). Moreover, every endomorphism \(f\) of I whose restriction to a nonzero submodule E of I is injective is injective (indeed, since (i) \(\Rightarrow\) (iii), I is an injective envelope of E, hence “E ∩ Ker \(f = 0\) ” implies “Ker \(f = 0\) ”. Given that, let \(f\) be a non-invertible element of \(\operatorname{End}_{\mathbf{A}}(\mathbf{M})\) ; by VIII, § 1, no. 4, prop. 9, it suffices to prove that \(1 - f\) is invertible. Since \(f\) is not injective, we have \(\operatorname{Ker} f \neq 0\) as the restriction of \(1 - f\) to \(\operatorname{Ker} f\) is injective, \(1 - f\) is injective, hence bijective.
\end{enumerate}

COROLLARY 1. --- The relation “I is a class of indecomposable injective A-modules” is collectivizing.

Indeed, by (iii) every indecomposable injective A-module is the injective envelope of a monogenic A-module.

COROLLARY 2. --- Let M be an A-module, and let I be an injective envelope of M. For I to be indecomposable, it is necessary and sufficient that 0 not be the intersection of two nonzero submodules of M.

The condition is necessary by Prop. 14 ((i) \(\Rightarrow\) (ii)). Conversely, if I is the direct sum of nonzero submodules \(\mathbf{I}_1\) and \(\mathbf{I}_2\) , we have:

\[
\mathrm{I} _ {1} \cap \mathrm{M} \neq 0, \quad \mathrm{I} _ {2} \cap \mathrm{M} \neq 0 \quad \text { and } \quad (\mathrm{I} _ {1} \cap \mathrm{M}) \cap (\mathrm{I} _ {2} \cap \mathrm{M}) = 0.
\]

Example 2. --- If A is commutative and Noetherian, the indecomposable injective A-modules are exactly the injective envelopes of the modules A/p where p is a prime ideal (X, p.~171, exercise 27).

